\haddockmoduleheading{System.Exit}
\label{module:System.Exit}
\haddockbeginheader
{\haddockverb\begin{verbatim}
module System.Exit (
    ExitCode(ExitSuccess, ExitFailure), exitWith, exitFailure, exitSuccess
  ) where\end{verbatim}}
\haddockendheader

\begin{haddockdesc}
\item[\begin{tabular}{@{}l}
data ExitCode
\end{tabular}]
{\haddockbegindoc
Defines the exit codes that a program can return.\par
\enspace \emph{Constructors}\par
\haddockbeginconstrs
\haddockdecltt{=} & \haddockdecltt{ExitSuccess} & indicates successful termination; \\
\haddockdecltt{|} & \haddockdecltt{ExitFailure Int} & indicates program failure with an exit code.
 The exact interpretation of the code is
 operating-system dependent.  In particular, some values
 may be prohibited (e.g. 0 on a POSIX-compliant system). \\
\end{tabulary}\par}
\end{haddockdesc}
\begin{haddockdesc}
\item[\begin{tabular}{@{}l}
instance Read ExitCode\\instance Show ExitCode\\instance Eq ExitCode\\instance Ord ExitCode\\
\end{tabular}]
\end{haddockdesc}
\begin{haddockdesc}
\item[\begin{tabular}{@{}l}
exitWith :: ExitCode -> IO a
\end{tabular}]
{\haddockbegindoc
Computation \haddockid{exitWith} \haddocktt{code} throws \haddockid{ExitCode} \haddocktt{code}.
 Normally this terminates the program, returning \haddocktt{code} to the
 program's caller.\par
On program termination, the standard \haddockid{Handle}s \haddockid{stdout} and
 \haddockid{stderr} are flushed automatically; any other buffered \haddockid{Handle}s
 need to be flushed manually, otherwise the buffered data will be
 discarded.\par
A program that fails in any other way is treated as if it had
 called \haddockid{exitFailure}.
 A program that terminates successfully without calling \haddockid{exitWith}
 explicitly is treated as if it had called \haddockid{exitWith} \haddockid{ExitSuccess}.\par
As an \haddockid{ExitCode} is an \haddockid{Exception}, it can be
 caught using the functions of \haddocktt{Control.Exception}.  This means that
 cleanup computations added with \haddockid{bracket} (from
 \haddocktt{Control.Exception}) are also executed properly on \haddockid{exitWith}.\par
Note: in GHC, \haddockid{exitWith} should be called from the main program
 thread in order to exit the process.  When called from another
 thread, \haddockid{exitWith} will throw an \haddockid{ExitCode} as normal, but the
 exception will not cause the process itself to exit.\par}
\end{haddockdesc}
\begin{haddockdesc}
\item[\begin{tabular}{@{}l}
exitFailure :: IO a
\end{tabular}]
{\haddockbegindoc
The computation \haddockid{exitFailure} is equivalent to
 \haddockid{exitWith} \haddocktt{(}\haddockid{ExitFailure} \emph{exitfail}\haddocktt{)},
 where \emph{exitfail} is implementation-dependent.\par}
\end{haddockdesc}
\begin{haddockdesc}
\item[\begin{tabular}{@{}l}
exitSuccess :: IO a
\end{tabular}]
{\haddockbegindoc
The computation \haddockid{exitSuccess} is equivalent to
 \haddockid{exitWith} \haddockid{ExitSuccess}, It terminates the program
 successfully.\par}
\end{haddockdesc}